\documentclass[12pt]{article}
\usepackage[utf8]{inputenc}
\usepackage[T1]{fontenc}
\usepackage{hyperref}
\title{Loxicha Zapotec Language Resources for K-12 Curriculum Development}
\author{Compiled by Claude Sonnet 5.5 (claude-sonnet-5-5) for Andr\'as Kornai}
\date{}
\begin{document}
\maketitle

\section{Introduction}
Loxicha Zapotec (ISO 639-3 ztp, Glottocode loxi1235; also called Western Pochutla Zapotec, endonym Diste) is an Otomanguean language of Oaxaca, Mexico. Wikitongues rates it EGIDS 6a (Vigorous). San Balt\'azar Loxicha Zapotec, which includes Santa Catarina Loxicha, is a distinct language and should not be confused with it. The searches found a few documentation resources, but no school materials.

\section{Primary Educational Resources}
No primers, textbooks or school programmes for ztp were found.

\section{Linguistic and Reference Materials}
\begin{itemize}
\item Wikipedia overview, whose references include Mikko Salminen's 2014 book \emph{D\'i'zt\`e': o zapoteco de San Agust\'in Loxicha, Oaxaca, M\'exico} (LINCOM, Munich): \url{https://en.wikipedia.org/wiki/Loxicha_Zapotec}
\item Glottolog entry: \url{https://glottolog.org/resource/languoid/id/loxi1235}
\item Wikitongues language page: \url{https://wikitongues.org/languages/ztp/}
\item ELDP documentation collection on the San Bartolom\'e Loxicha variety (the page says the recordings were used for a grammar and dictionary): \url{https://eldp.access.preservica.com/dk01620182/}
\item ASJP word list: \url{https://asjp.clld.org/languages/ZAPOTEC_LOXICHA}
\item PHOIBLE phonological inventory: \url{https://phoible.org/languages/loxi1235}
\end{itemize}
A talking dictionary exists for a nearby variety, which the source calls Miahuatec Zapotec (San Bartolom\'e Loxicha); that is a different variety, so check its fit: \url{https://works.swarthmore.edu/fac-linguistics/361}

\section{Religious and Community Resources}
None were found.

\section{Digital and Audio Resources}
\begin{itemize}
\item Global Recordings Network page: \url{https://glorecnet-website.k.globalrecordings.net/en/language/ztp}
\end{itemize}

\section{Curriculum Development Context}
Loxicha Zapotec is still vigorous and has a documentation base: the ELDP collection and the Salminen book. No school-oriented materials were found. Curriculum development could build on the documentation, with an agreed orthography and community teachers. The related dictionary and the documentation archives are worth checking for reusable vocabulary and texts.
\end{document}